% ============================================================
% 第 10 课 · LaTeX 最小练习（先编出你的第一个 PDF）
%
% 编译（中文必须用 xelatex；编两遍目录/引用才准）：
%     cd ~/tjurm2027-test1/report
%     xelatex lesson10_drill.tex
%
% 目标：把 3 个 TODO 填上，编出一个含「章节 / 公式 / 代码块 / 图片」的 PDF。
% 提示：LaTeX 里 % 之后是注释；\ 开头的都是命令；{} 里是参数。
% ============================================================
\documentclass[11pt,a4paper]{ctexart}   % 中文文档用 ctexart + xelatex

\usepackage{graphicx}    % 插图
\usepackage{listings}    % 代码块
\usepackage{amsmath}     % 数学公式
\usepackage{xcolor}      % 代码块高亮

\title{我的第一个 \LaTeX{} 文档}
\author{Liboven}
\date{\today}

\begin{document}
\maketitle

\section{一句话介绍}
这是用 \LaTeX{} 排出来的：中文正常显示，章节编号自动生成。
注意：正文里的下划线、百分号等特殊字符要转义，比如 \texttt{my\_strlen}、50\%。

% ------------------------------------------------------------------
% TODO 1：加一个二级标题「数学公式」，并在下面放一个带编号的公式
   \subsection{数学公式}
   \begin{equation}
     V = 0.2989R + 0.5870G + 0.1140B
   \end{equation}
% ------------------------------------------------------------------


% ------------------------------------------------------------------
% TODO 2：加一个二级标题「代码块」，并在下面写三行 C++ 代码
   \subsection{代码块}
   \begin{lstlisting}
   while (*p != '\0') { ++p; }
   \end{lstlisting}
%   注意：代码块里的内容不用转义
% ------------------------------------------------------------------


% ------------------------------------------------------------------
% TODO 3：加一个二级标题「图片」，插入仓库里的一张结果图
   \subsection{图片}
   \begin{figure}[h]
     \centering
     \includegraphics[width=0.6\textwidth]{../images/rgb2gray/output.jpg}
     \caption{我的 rgb2gray 输出}
   \end{figure}
%   图片路径是相对这个 .tex 文件的
% ------------------------------------------------------------------

\end{document}
